% =============================================================================
% example-technical-user-manual.tex
% Style: latex-docs-technical-user-manual
% Build: pdflatex via latexmk; minted-aware with listings fallback.
% =============================================================================
\documentclass[11pt]{article}
\usepackage{latex-docs-technical-user-manual}

\setDocTitle{AetherForge Engine -- Quick Start Guide}
\setDocSubtitle{Your first scene in under five minutes}
\setDocVersion{v1.0}
\setDocStatus{Released}
\setDocOwner{Documentation}
\setDocAuthor{Jordan Suber}
\setDocDate{May 8, 2026}
\setDocSummary{%
  Reference user manual demonstrating the user-manual subtype: tasks
  with stepped procedures, GUI-element typography, and a
  troubleshooting Q/A. Every paragraph is tuned for an end-user reading
  level.%
}

\begin{document}
\makedoctocpage

\section{Before you begin}

This guide walks you through opening AetherForge for the first time,
loading the sample scene, and rendering a single frame. You should be
able to complete it in under five minutes on a system that already has
the engine installed.

\begin{UserNote}[What you'll need]
A working installation of AetherForge Engine, and a sample scene file
(included with the installer at \filename{samples/hello.afs}).
\end{UserNote}

\section{Tasks}

\begin{task}{Open AetherForge for the first time}{%
  Get the application running and confirm the welcome screen appears.}
\begin{stepresult}
  \SR{Launch AetherForge from your applications menu.}%
     {The \Window{Welcome to AetherForge} window appears with three
      large buttons: \Button{Open scene}, \Button{New scene}, and
      \Button{Tutorials}.}
  \SR{Click \Button{Tutorials}.}%
     {The tutorial library opens in the main editor window.}
  \SR{Close the tutorial library by clicking \Menu{File > Close} or
      pressing \kbd{Ctrl+W}.}%
     {You are returned to the empty editor.}
\end{stepresult}
\end{task}

\begin{task}{Load the sample scene}{%
  Open the bundled \texttt{hello.afs} sample.}
\begin{stepresult}
  \SR{Click \Menu{File > Open Scene}.}%
     {The \Window{Open Scene} dialog appears.}
  \SR{In the \Field{Path} field, enter the full path to
      \filename{samples/hello.afs} (or use the \Button{Browse} button).}%
     {The path is shown; the \Button{Open} button becomes active.}
  \SR{Click \Button{Open}.}%
     {The sample scene loads. You should see a single textured cube in
      the viewport.}
\end{stepresult}
\end{task}

\begin{task}{Render a single frame}{%
  Trigger an offline render and save the output as PNG.}
\begin{stepresult}
  \SR{Click \Menu{Render > Render Single Frame}.}%
     {A progress bar appears in the status area at the bottom of the
      editor.}
  \SR{When the render completes, the \Window{Save rendered image}
      dialog appears. In the \Field{Filename} field, type
      \filename{first-frame.png}.}%
     {The filename is accepted; the \Button{Save} button becomes
      active.}
  \SR{Click \Button{Save}.}%
     {The PNG is written to the chosen location. The status area shows
      \emph{Render saved}.}
\end{stepresult}
\end{task}

\begin{UserTip}[Re-run faster]
You can re-render the current scene with \kbd{F12} once you have used
\Menu{Render > Render Single Frame} at least once in this session.
\end{UserTip}

\begin{troubleshooting}
  \TS{The welcome window does not appear when I launch AetherForge.}%
     {Check that no other instance is already running. On Linux, look
      for a stale lockfile at \filename{\$HOME/.aetherforge/lockfile}.}
  \TS{The sample scene loads but the cube is missing or invisible.}%
     {Press \kbd{Home} to recentre the view on the scene's origin. If
      the cube is still missing, the scene file may be corrupted; copy
      a fresh one from the installer's \filename{samples/} folder.}
  \TS{The render completes but the saved PNG is all black.}%
     {Confirm at least one light source is present in the scene
      (\Menu{Scene > Lighting > Show light list}) and that GPU rendering
      is enabled in \Menu{Edit > Preferences > Rendering}.}
\end{troubleshooting}

\section{Where to next}

\begin{itemize}
  \item Read the design specification for the engine's render pipeline.
  \item Try the \emph{Animating a camera} tutorial, which builds on the
        scene you just loaded.
  \item File issues at \texttt{aetherforge/issues} on the project
        repository.
\end{itemize}

\end{document}
